\documentclass[12pt]{article}
\usepackage{amsmath,amssymb}

\begin{document}
\section*{{2020 AIME II Problems/Problem 15}}
Source: AoPS Wiki

\subsection*{Solution}
Let $O$ be the circumcenter of $\triangle ABC$ ; say $OT$ intersects $BC$ at $M$ ; draw segments $XM$ , and $YM$ . We have $MT=3\sqrt{15}$ . Since $\angle A=\angle CBT=\angle BCT$ , we have $\cos A=\tfrac{11}{16}$ . Notice that $AXTY$ is cyclic, so $\angle XTY=180^{\circ}-A$ , so $\cos XTY=-\cos A$ , and the cosine law in $\triangle TXY$ gives \[1143-2XY^2=-\frac{11}{8}\cdot XT\cdot YT.\] Since $\triangle BMT \cong \triangle CMT$ , we have $TM\perp BC$ , and therefore quadrilaterals $BXTM$ and $CYTM$ are cyclic. Let $P$ (resp. $Q$ ) be the midpoint of $BT$ (resp. $CT$ ). So $P$ (resp. $Q$ ) is the center of $(BXTM)$ (resp. $CYTM$ ). Then $\theta=\angle ABC=\angle MTX$ and $\phi=\angle ACB=\angle YTM$ . So $\angle XPM=2\theta$ , so \[\frac{\frac{XM}{2}}{XP}=\sin \theta,\] which yields $XM=2XP\sin \theta=BT(=CT)\sin \theta=TY$ . Similarly we have $YM=XT$ . Ptolemy's theorem in $BXTM$ gives \[16TY=11TX+3\sqrt{15}BX,\] while Pythagoras' theorem gives $BX^2+XT^2=16^2$ . Similarly, Ptolemy's theorem in $YTMC$ gives \[16TX=11TY+3\sqrt{15}CY\] while Pythagoras' theorem in $\triangle CYT$ gives $CY^2+YT^2=16^2$ . Solve this for $XT$ and $TY$ and substitute into the equation about $\cos XTY$ to obtain the result $XY^2=\boxed{717}$ . (Notice that $MXTY$ is a parallelogram, which is an important theorem in Olympiad, and there are some other ways of computation under this observation.) -Fanyuchen20020715

\end{document}
